% --- Archon physics formalization source begin ---
% archon:physics
% archon:covers PhyXMiniProblems/problem_phyx_mini_0435.lean
% archon:source-report reports/phyx_mini/problem_phyx_mini_0435.source.json

\chapter{Physics problem phyx\_mini\_0435}
\label{ch:PhyXMiniProblems_problem_phyx_mini_0435}

\paragraph{Problem source.}
\#\# Physical scenario

Figure shows a \$pV\$ diagram for a thermodynamic cycle, with the device now operated in reverse.

\#\# Question

Determine the values of \$Q\_H\$.

\#\# Answer choices

- A: A: 22.28 \$\textbackslash{}times 10\textasciicircum{}4 \textbackslash{}

- B: B: 26.33 \$\textbackslash{}times 10\textasciicircum{}5 \textbackslash{}

- C: C: 15.19 \$\textbackslash{}times 10\textasciicircum{}5 \textbackslash{}

- D: D: 22.28 \$\textbackslash{}times 10\textasciicircum{}5 \textbackslash{}

\#\# Figure caption (auxiliary; use the image as primary evidence)

The image is a pressure-volume (p-V) diagram with a cyclic thermodynamic process represented by a square path. 

- **Axes:** The vertical axis is labeled \textbackslash{}( p \textbackslash{}) (atm) for pressure in atmospheres, and the horizontal axis is labeled \textbackslash{}( V \textbackslash{}) (m\textbackslash{}(\textasciicircum{}3\textbackslash{})) for volume in cubic meters.
- **Path:** The process goes through four states, numbered 1 to 4, forming a square:
  - State 1 to 2: Moves horizontally to the right.
  - State 2 to 3: Moves vertically downwards.
  - State 3 to 4: Moves horizontally to the left.
  - State 4 to 1: Moves vertically upwards.
- **Arrows:** Indicate the direction of the process flow, counter-clockwise around the square.
- **Temperature Lines:** Two red dashed lines are shown:
  - A line passing through states 1 and 2 is labeled "2700 K".
  - Another line passing through states 3 and 4 is labeled "300 K".

The diagram visually represents a thermodynamic cycle with specified temperatures.

\#\# Dataset metadata

Laws of Thermodynamics; Physical Model Grounding Reasoning, Multi-Formula Reasoning

\paragraph{Recorded answer/context.}
D: D: 22.28 \$\textbackslash{}times 10\textasciicircum{}5 \textbackslash{}

\paragraph{Figure/image path.}
/root/proposal\_for\_physic/science-mango/phyx\_mini\_run/phyx\_data/test\_image/435.png

\paragraph{Formalization target.}
create a compiling Lean file with sorry bodies at `PhyXMiniProblems/problem\_phyx\_mini\_0435.lean`. The Lean declarations must preserve the physical quantities, dimensions or dimensional roles, figure labels, governing-law hypotheses, and final relation expressed by this problem.
Use Mathlib/Physlib names found through LeanExplore where available. If a physics API is missing, introduce faithful local abstractions rather than scalar placeholder aliases.

\begin{theorem}[Physics formalization target]
\label{thm:physics:phyx_mini_0435:target}
\lean{PhyXMiniProblems.ProblemPhyXMini0435.hotReservoirHeat_matches_recordedAnswerD}
\uses{def:physics:phyx-mini-0435:phyxminiproblems-problemphyxmini0435-reversedrectangularidealgascycle, def:physics:phyx-mini-0435:phyxminiproblems-problemphyxmini0435-matchesproblemandprimaryfigure, def:physics:phyx-mini-0435:phyxminiproblems-problemphyxmini0435-usesidealgascvoverr2model, def:physics:phyx-mini-0435:phyxminiproblems-problemphyxmini0435-hasphysicalcycleparameters, def:physics:phyx-mini-0435:phyxminiproblems-problemphyxmini0435-satisfiesidealgaslaws, def:physics:phyx-mini-0435:phyxminiproblems-problemphyxmini0435-hotreservoirheatinjoules, def:physics:phyx-mini-0435:phyxminiproblems-problemphyxmini0435-recordedanswerchoice, def:physics:phyx-mini-0435:phyxminiproblems-problemphyxmini0435-compatiblewithdisplayedheat, def:physics:phyx-mini-0435:phyxminiproblems-problemphyxmini0435-isuniqueclosestdisplayedheat}
The assigned autoformalize agent should translate this physics problem into Lean declarations in the covered file, with theorem and lemma proof bodies written as `by sorry`.
\end{theorem}
\begin{proof}
This is an autoformalization task, not a proof task. Produce statements that can later be proved by the physics prover without weakening the physical model.
\end{proof}

% --- Archon declaration topology begin ---
\section{Declaration topology}
\begin{definition}[Volume Quantity]
  \label{def:physics:phyx-mini-0435:phyxminiproblems-problemphyxmini0435-volumequantity}
  \lean{PhyXMiniProblems.ProblemPhyXMini0435.VolumeQuantity}
  A nonnegative physical volume carrying dimension L³
\end{definition}
\begin{proof}
  This is the declared model definition of Volume Quantity.
\end{proof}
\begin{definition}[volume In Cubic Meters]
  \label{def:physics:phyx-mini-0435:phyxminiproblems-problemphyxmini0435-volumeincubicmeters}
  \lean{PhyXMiniProblems.ProblemPhyXMini0435.volumeInCubicMeters}
  \uses{def:physics:phyx-mini-0435:phyxminiproblems-problemphyxmini0435-volumequantity}
  Cubic-metre readout of a physical volume
\end{definition}
\begin{proof}
  This is the declared model definition of volume In Cubic Meters.
\end{proof}
\begin{definition}[pressure In Pascals]
  \label{def:physics:phyx-mini-0435:phyxminiproblems-problemphyxmini0435-pressureinpascals}
  \lean{PhyXMiniProblems.ProblemPhyXMini0435.pressureInPascals}
  Pascal readout of a physical pressure
\end{definition}
\begin{proof}
  This is the declared model definition of pressure In Pascals.
\end{proof}
\begin{definition}[pressure In Atmospheres]
  \label{def:physics:phyx-mini-0435:phyxminiproblems-problemphyxmini0435-pressureinatmospheres}
  \lean{PhyXMiniProblems.ProblemPhyXMini0435.pressureInAtmospheres}
  Standard-atmosphere readout of a physical pressure
\end{definition}
\begin{proof}
  This is the declared model definition of pressure In Atmospheres.
\end{proof}
\begin{definition}[energy In Joules]
  \label{def:physics:phyx-mini-0435:phyxminiproblems-problemphyxmini0435-energyinjoules}
  \lean{PhyXMiniProblems.ProblemPhyXMini0435.energyInJoules}
  Joule readout of a signed physical energy, heat transfer, or work
\end{definition}
\begin{proof}
  This is the declared model definition of energy In Joules.
\end{proof}
\begin{definition}[temperature In Kelvin]
  \label{def:physics:phyx-mini-0435:phyxminiproblems-problemphyxmini0435-temperatureinkelvin}
  \lean{PhyXMiniProblems.ProblemPhyXMini0435.temperatureInKelvin}
  Kelvin readout of a Physlib absolute temperature stored in an explicitly chosen zero-preserving temperature unit
\end{definition}
\begin{proof}
  This is the declared model definition of temperature In Kelvin.
\end{proof}
\begin{lemma}[standard Atmosphere In Pascals]
  \label{lem:physics:phyx-mini-0435:phyxminiproblems-problemphyxmini0435-standardatmosphereinpascals}
  \lean{PhyXMiniProblems.ProblemPhyXMini0435.standardAtmosphereInPascals}
  \uses{def:physics:phyx-mini-0435:phyxminiproblems-problemphyxmini0435-pressureinpascals}
  Physlib's standard-atmosphere unit has the exact SI value 101325 Pa
\end{lemma}
\begin{proof}
  The conclusion follows from the governing relations and figure readouts named in the statement of standard Atmosphere In Pascals.
\end{proof}
\begin{definition}[Cycle State]
  \label{def:physics:phyx-mini-0435:phyxminiproblems-problemphyxmini0435-cyclestate}
  \lean{PhyXMiniProblems.ProblemPhyXMini0435.CycleState}
  The four numbered states printed on the rectangular cycle
\end{definition}
\begin{proof}
  This is the declared model definition of Cycle State.
\end{proof}
\begin{definition}[Cycle Leg]
  \label{def:physics:phyx-mini-0435:phyxminiproblems-problemphyxmini0435-cycleleg}
  \lean{PhyXMiniProblems.ProblemPhyXMini0435.CycleLeg}
  The four directed legs, in the counter-clockwise order shown
\end{definition}
\begin{proof}
  This is the declared model definition of Cycle Leg.
\end{proof}
\begin{definition}[leg Source]
  \label{def:physics:phyx-mini-0435:phyxminiproblems-problemphyxmini0435-legsource}
  \lean{PhyXMiniProblems.ProblemPhyXMini0435.legSource}
  \uses{def:physics:phyx-mini-0435:phyxminiproblems-problemphyxmini0435-cyclestate, def:physics:phyx-mini-0435:phyxminiproblems-problemphyxmini0435-cycleleg}
  Initial state of a directed reverse-cycle leg
\end{definition}
\begin{proof}
  This is the declared model definition of leg Source.
\end{proof}
\begin{definition}[leg Target]
  \label{def:physics:phyx-mini-0435:phyxminiproblems-problemphyxmini0435-legtarget}
  \lean{PhyXMiniProblems.ProblemPhyXMini0435.legTarget}
  \uses{def:physics:phyx-mini-0435:phyxminiproblems-problemphyxmini0435-cyclestate, def:physics:phyx-mini-0435:phyxminiproblems-problemphyxmini0435-cycleleg}
  Final state of a directed reverse-cycle leg
\end{definition}
\begin{proof}
  This is the declared model definition of leg Target.
\end{proof}
\begin{definition}[Process Kind]
  \label{def:physics:phyx-mini-0435:phyxminiproblems-problemphyxmini0435-processkind}
  \lean{PhyXMiniProblems.ProblemPhyXMini0435.ProcessKind}
  Thermodynamic constraint represented by a side of the rectangle
\end{definition}
\begin{proof}
  This is the declared model definition of Process Kind.
\end{proof}
\begin{definition}[displayed Process Kind]
  \label{def:physics:phyx-mini-0435:phyxminiproblems-problemphyxmini0435-displayedprocesskind}
  \lean{PhyXMiniProblems.ProblemPhyXMini0435.displayedProcessKind}
  \uses{def:physics:phyx-mini-0435:phyxminiproblems-problemphyxmini0435-cycleleg, def:physics:phyx-mini-0435:phyxminiproblems-problemphyxmini0435-processkind}
  Process kind read from the horizontal or vertical geometry of each leg
\end{definition}
\begin{proof}
  This is the declared model definition of displayed Process Kind.
\end{proof}
\begin{definition}[Figure Axis]
  \label{def:physics:phyx-mini-0435:phyxminiproblems-problemphyxmini0435-figureaxis}
  \lean{PhyXMiniProblems.ProblemPhyXMini0435.FigureAxis}
  The two Cartesian axes in the supplied diagram
\end{definition}
\begin{proof}
  This is the declared model definition of Figure Axis.
\end{proof}
\begin{definition}[Axis Quantity]
  \label{def:physics:phyx-mini-0435:phyxminiproblems-problemphyxmini0435-axisquantity}
  \lean{PhyXMiniProblems.ProblemPhyXMini0435.AxisQuantity}
  Physical quantity assigned to a diagram axis
\end{definition}
\begin{proof}
  This is the declared model definition of Axis Quantity.
\end{proof}
\begin{definition}[Axis Display Unit]
  \label{def:physics:phyx-mini-0435:phyxminiproblems-problemphyxmini0435-axisdisplayunit}
  \lean{PhyXMiniProblems.ProblemPhyXMini0435.AxisDisplayUnit}
  Units printed next to the axes in the bitmap
\end{definition}
\begin{proof}
  This is the declared model definition of Axis Display Unit.
\end{proof}
\begin{definition}[Arrow Direction]
  \label{def:physics:phyx-mini-0435:phyxminiproblems-problemphyxmini0435-arrowdirection}
  \lean{PhyXMiniProblems.ProblemPhyXMini0435.ArrowDirection}
  Direction of a visible arrowhead on one side of the rectangle
\end{definition}
\begin{proof}
  This is the declared model definition of Arrow Direction.
\end{proof}
\begin{definition}[displayed Arrow Direction]
  \label{def:physics:phyx-mini-0435:phyxminiproblems-problemphyxmini0435-displayedarrowdirection}
  \lean{PhyXMiniProblems.ProblemPhyXMini0435.displayedArrowDirection}
  \uses{def:physics:phyx-mini-0435:phyxminiproblems-problemphyxmini0435-cycleleg, def:physics:phyx-mini-0435:phyxminiproblems-problemphyxmini0435-arrowdirection}
  Arrow direction visible on each reverse-cycle leg
\end{definition}
\begin{proof}
  This is the declared model definition of displayed Arrow Direction.
\end{proof}
\begin{definition}[Cycle Orientation]
  \label{def:physics:phyx-mini-0435:phyxminiproblems-problemphyxmini0435-cycleorientation}
  \lean{PhyXMiniProblems.ProblemPhyXMini0435.CycleOrientation}
  Orientation of the closed path in the pV plane
\end{definition}
\begin{proof}
  This is the declared model definition of Cycle Orientation.
\end{proof}
\begin{definition}[Isotherm Label]
  \label{def:physics:phyx-mini-0435:phyxminiproblems-problemphyxmini0435-isothermlabel}
  \lean{PhyXMiniProblems.ProblemPhyXMini0435.IsothermLabel}
  The two dashed temperature annotations in the primary bitmap
\end{definition}
\begin{proof}
  This is the declared model definition of Isotherm Label.
\end{proof}
\begin{definition}[Pressure Volume Cycle Figure]
  \label{def:physics:phyx-mini-0435:phyxminiproblems-problemphyxmini0435-pressurevolumecyclefigure}
  \lean{PhyXMiniProblems.ProblemPhyXMini0435.PressureVolumeCycleFigure}
  \uses{def:physics:phyx-mini-0435:phyxminiproblems-problemphyxmini0435-volumequantity, def:physics:phyx-mini-0435:phyxminiproblems-problemphyxmini0435-cyclestate, def:physics:phyx-mini-0435:phyxminiproblems-problemphyxmini0435-cycleleg, def:physics:phyx-mini-0435:phyxminiproblems-problemphyxmini0435-figureaxis, def:physics:phyx-mini-0435:phyxminiproblems-problemphyxmini0435-axisquantity, def:physics:phyx-mini-0435:phyxminiproblems-problemphyxmini0435-axisdisplayunit, def:physics:phyx-mini-0435:phyxminiproblems-problemphyxmini0435-arrowdirection, def:physics:phyx-mini-0435:phyxminiproblems-problemphyxmini0435-isothermlabel}
  Literal geometric and physical content of the supplied pV diagram
\end{definition}
\begin{proof}
  This is the declared model definition of Pressure Volume Cycle Figure.
\end{proof}
\begin{definition}[Thermodynamic Device Role]
  \label{def:physics:phyx-mini-0435:phyxminiproblems-problemphyxmini0435-thermodynamicdevicerole}
  \lean{PhyXMiniProblems.ProblemPhyXMini0435.ThermodynamicDeviceRole}
  Physical role of the device after the stated reversal
\end{definition}
\begin{proof}
  This is the declared model definition of Thermodynamic Device Role.
\end{proof}
\begin{definition}[Gas Model]
  \label{def:physics:phyx-mini-0435:phyxminiproblems-problemphyxmini0435-gasmodel}
  \lean{PhyXMiniProblems.ProblemPhyXMini0435.GasModel}
  Constitutive class of the working gas
\end{definition}
\begin{proof}
  This is the declared model definition of Gas Model.
\end{proof}
\begin{definition}[Reversed Rectangular Ideal Gas Cycle]
  \label{def:physics:phyx-mini-0435:phyxminiproblems-problemphyxmini0435-reversedrectangularidealgascycle}
  \lean{PhyXMiniProblems.ProblemPhyXMini0435.ReversedRectangularIdealGasCycle}
  \uses{def:physics:phyx-mini-0435:phyxminiproblems-problemphyxmini0435-cyclestate, def:physics:phyx-mini-0435:phyxminiproblems-problemphyxmini0435-cycleleg, def:physics:phyx-mini-0435:phyxminiproblems-problemphyxmini0435-cycleorientation, def:physics:phyx-mini-0435:phyxminiproblems-problemphyxmini0435-pressurevolumecyclefigure, def:physics:phyx-mini-0435:phyxminiproblems-problemphyxmini0435-thermodynamicdevicerole, def:physics:phyx-mini-0435:phyxminiproblems-problemphyxmini0435-gasmodel}
  The reverse-cycle device and its physical observables. Heat is positive into the gas and boundary work is positive when done by the gas. The scalar sampleGasConstantJoulesPerKelvin is the SI readout of nR for the complete gas sample, whose amount is not supplied separately in the source. The constant-volume heat capacity is retained as a separate SI readout because the source does not state it, while the recorded numerical answer depends on its value
\end{definition}
\begin{proof}
  This is the declared model definition of Reversed Rectangular Ideal Gas Cycle.
\end{proof}
\begin{definition}[gas Temperature In Kelvin]
  \label{def:physics:phyx-mini-0435:phyxminiproblems-problemphyxmini0435-gastemperatureinkelvin}
  \lean{PhyXMiniProblems.ProblemPhyXMini0435.gasTemperatureInKelvin}
  \uses{def:physics:phyx-mini-0435:phyxminiproblems-problemphyxmini0435-temperatureinkelvin, def:physics:phyx-mini-0435:phyxminiproblems-problemphyxmini0435-cyclestate, def:physics:phyx-mini-0435:phyxminiproblems-problemphyxmini0435-reversedrectangularidealgascycle}
  Kelvin readout at a numbered state
\end{definition}
\begin{proof}
  This is the declared model definition of gas Temperature In Kelvin.
\end{proof}
\begin{definition}[internal Energy In Joules]
  \label{def:physics:phyx-mini-0435:phyxminiproblems-problemphyxmini0435-internalenergyinjoules}
  \lean{PhyXMiniProblems.ProblemPhyXMini0435.internalEnergyInJoules}
  \uses{def:physics:phyx-mini-0435:phyxminiproblems-problemphyxmini0435-energyinjoules, def:physics:phyx-mini-0435:phyxminiproblems-problemphyxmini0435-cyclestate, def:physics:phyx-mini-0435:phyxminiproblems-problemphyxmini0435-reversedrectangularidealgascycle}
  Joule readout of internal energy at a numbered state
\end{definition}
\begin{proof}
  This is the declared model definition of internal Energy In Joules.
\end{proof}
\begin{definition}[work Done By Gas In Joules]
  \label{def:physics:phyx-mini-0435:phyxminiproblems-problemphyxmini0435-workdonebygasinjoules}
  \lean{PhyXMiniProblems.ProblemPhyXMini0435.workDoneByGasInJoules}
  \uses{def:physics:phyx-mini-0435:phyxminiproblems-problemphyxmini0435-energyinjoules, def:physics:phyx-mini-0435:phyxminiproblems-problemphyxmini0435-cycleleg, def:physics:phyx-mini-0435:phyxminiproblems-problemphyxmini0435-reversedrectangularidealgascycle}
  Signed joule readout of boundary work done by the gas on a leg
\end{definition}
\begin{proof}
  This is the declared model definition of work Done By Gas In Joules.
\end{proof}
\begin{definition}[heat Transferred Into Gas In Joules]
  \label{def:physics:phyx-mini-0435:phyxminiproblems-problemphyxmini0435-heattransferredintogasinjoules}
  \lean{PhyXMiniProblems.ProblemPhyXMini0435.heatTransferredIntoGasInJoules}
  \uses{def:physics:phyx-mini-0435:phyxminiproblems-problemphyxmini0435-energyinjoules, def:physics:phyx-mini-0435:phyxminiproblems-problemphyxmini0435-cycleleg, def:physics:phyx-mini-0435:phyxminiproblems-problemphyxmini0435-reversedrectangularidealgascycle}
  Signed joule readout of heat transferred into the gas on a leg
\end{definition}
\begin{proof}
  This is the declared model definition of heat Transferred Into Gas In Joules.
\end{proof}
\begin{definition}[Matches Problem And Primary Figure]
  \label{def:physics:phyx-mini-0435:phyxminiproblems-problemphyxmini0435-matchesproblemandprimaryfigure}
  \lean{PhyXMiniProblems.ProblemPhyXMini0435.MatchesProblemAndPrimaryFigure}
  \uses{def:physics:phyx-mini-0435:phyxminiproblems-problemphyxmini0435-volumeincubicmeters, def:physics:phyx-mini-0435:phyxminiproblems-problemphyxmini0435-pressureinatmospheres, def:physics:phyx-mini-0435:phyxminiproblems-problemphyxmini0435-displayedarrowdirection, def:physics:phyx-mini-0435:phyxminiproblems-problemphyxmini0435-reversedrectangularidealgascycle, def:physics:phyx-mini-0435:phyxminiproblems-problemphyxmini0435-gastemperatureinkelvin}
  Exact transcription of the prose and primary bitmap. The dashed 2700 K curve visibly meets state 2 and the dashed 300 K curve visibly meets state 4; the auxiliary caption's claim that each passes through two vertices is not used because the bitmap is designated as primary evidence
\end{definition}
\begin{proof}
  This is the declared model definition of Matches Problem And Primary Figure.
\end{proof}
\begin{definition}[Uses Ideal Gas Cv Over R2 Model]
  \label{def:physics:phyx-mini-0435:phyxminiproblems-problemphyxmini0435-usesidealgascvoverr2model}
  \lean{PhyXMiniProblems.ProblemPhyXMini0435.UsesIdealGasCvOverR2Model}
  \uses{def:physics:phyx-mini-0435:phyxminiproblems-problemphyxmini0435-reversedrectangularidealgascycle}
  The supplied source omits the heat capacity of the working gas. The recorded answer D is the result for an ideal gas with C\_V/(nR) = 2; this predicate makes that necessary constitutive choice explicit instead of hiding it in a definition of Q\_H
\end{definition}
\begin{proof}
  This is the declared model definition of Uses Ideal Gas Cv Over R2 Model.
\end{proof}
\begin{definition}[Has Physical Cycle Parameters]
  \label{def:physics:phyx-mini-0435:phyxminiproblems-problemphyxmini0435-hasphysicalcycleparameters}
  \lean{PhyXMiniProblems.ProblemPhyXMini0435.HasPhysicalCycleParameters}
  \uses{def:physics:phyx-mini-0435:phyxminiproblems-problemphyxmini0435-volumeincubicmeters, def:physics:phyx-mini-0435:phyxminiproblems-problemphyxmini0435-pressureinpascals, def:physics:phyx-mini-0435:phyxminiproblems-problemphyxmini0435-reversedrectangularidealgascycle, def:physics:phyx-mini-0435:phyxminiproblems-problemphyxmini0435-gastemperatureinkelvin}
  Positivity and nondegeneracy conditions for a physical gas cycle
\end{definition}
\begin{proof}
  This is the declared model definition of Has Physical Cycle Parameters.
\end{proof}
\begin{definition}[Satisfies Ideal Gas Laws]
  \label{def:physics:phyx-mini-0435:phyxminiproblems-problemphyxmini0435-satisfiesidealgaslaws}
  \lean{PhyXMiniProblems.ProblemPhyXMini0435.SatisfiesIdealGasLaws}
  \uses{def:physics:phyx-mini-0435:phyxminiproblems-problemphyxmini0435-volumeincubicmeters, def:physics:phyx-mini-0435:phyxminiproblems-problemphyxmini0435-pressureinpascals, def:physics:phyx-mini-0435:phyxminiproblems-problemphyxmini0435-legsource, def:physics:phyx-mini-0435:phyxminiproblems-problemphyxmini0435-legtarget, def:physics:phyx-mini-0435:phyxminiproblems-problemphyxmini0435-displayedprocesskind, def:physics:phyx-mini-0435:phyxminiproblems-problemphyxmini0435-reversedrectangularidealgascycle, def:physics:phyx-mini-0435:phyxminiproblems-problemphyxmini0435-gastemperatureinkelvin, def:physics:phyx-mini-0435:phyxminiproblems-problemphyxmini0435-internalenergyinjoules, def:physics:phyx-mini-0435:phyxminiproblems-problemphyxmini0435-workdonebygasinjoules, def:physics:phyx-mini-0435:phyxminiproblems-problemphyxmini0435-heattransferredintogasinjoules}
  Macroscopic governing laws for the chosen ideal-gas model: pV = (nR)T at each equilibrium state; U = C\_V T for the separately modeled constant-volume heat capacity; zero boundary work on an isochoric leg; W\_by = p(V\_f - V\_i) on an isobaric leg; Q\_in = U\_f - U\_i + W\_by on every directed leg. No field specifies a signed leg heat, total hot-side heat, or answer choice
\end{definition}
\begin{proof}
  This is the declared model definition of Satisfies Ideal Gas Laws.
\end{proof}
\begin{definition}[hot Reservoir Heat In Joules]
  \label{def:physics:phyx-mini-0435:phyxminiproblems-problemphyxmini0435-hotreservoirheatinjoules}
  \lean{PhyXMiniProblems.ProblemPhyXMini0435.hotReservoirHeatInJoules}
  \uses{def:physics:phyx-mini-0435:phyxminiproblems-problemphyxmini0435-reversedrectangularidealgascycle, def:physics:phyx-mini-0435:phyxminiproblems-problemphyxmini0435-heattransferredintogasinjoules}
  Positive heat delivered by the gas to the hot reservoir. The two legs on which heat leaves the gas are 2 → 1 and 1 → 4, so the sign is reversed. This definition contains no numerical answer
\end{definition}
\begin{proof}
  This is the declared model definition of hot Reservoir Heat In Joules.
\end{proof}
\begin{definition}[cold Reservoir Heat In Joules]
  \label{def:physics:phyx-mini-0435:phyxminiproblems-problemphyxmini0435-coldreservoirheatinjoules}
  \lean{PhyXMiniProblems.ProblemPhyXMini0435.coldReservoirHeatInJoules}
  \uses{def:physics:phyx-mini-0435:phyxminiproblems-problemphyxmini0435-reversedrectangularidealgascycle, def:physics:phyx-mini-0435:phyxminiproblems-problemphyxmini0435-heattransferredintogasinjoules}
  Positive heat absorbed by the gas on the low-temperature side
\end{definition}
\begin{proof}
  This is the declared model definition of cold Reservoir Heat In Joules.
\end{proof}
\begin{definition}[cycle Work Input In Joules]
  \label{def:physics:phyx-mini-0435:phyxminiproblems-problemphyxmini0435-cycleworkinputinjoules}
  \lean{PhyXMiniProblems.ProblemPhyXMini0435.cycleWorkInputInJoules}
  \uses{def:physics:phyx-mini-0435:phyxminiproblems-problemphyxmini0435-reversedrectangularidealgascycle, def:physics:phyx-mini-0435:phyxminiproblems-problemphyxmini0435-workdonebygasinjoules}
  Positive work supplied to drive one traversal of the reverse cycle
\end{definition}
\begin{proof}
  This is the declared model definition of cycle Work Input In Joules.
\end{proof}
\begin{definition}[Answer Choice]
  \label{def:physics:phyx-mini-0435:phyxminiproblems-problemphyxmini0435-answerchoice}
  \lean{PhyXMiniProblems.ProblemPhyXMini0435.AnswerChoice}
  Labels of the four displayed answers
\end{definition}
\begin{proof}
  This is the declared model definition of Answer Choice.
\end{proof}
\begin{definition}[displayed Hot Reservoir Heat In Joules]
  \label{def:physics:phyx-mini-0435:phyxminiproblems-problemphyxmini0435-displayedhotreservoirheatinjoules}
  \lean{PhyXMiniProblems.ProblemPhyXMini0435.displayedHotReservoirHeatInJoules}
  \uses{def:physics:phyx-mini-0435:phyxminiproblems-problemphyxmini0435-answerchoice}
  Joule value printed beside each answer label
\end{definition}
\begin{proof}
  This is the declared model definition of displayed Hot Reservoir Heat In Joules.
\end{proof}
\begin{definition}[recorded Answer Choice]
  \label{def:physics:phyx-mini-0435:phyxminiproblems-problemphyxmini0435-recordedanswerchoice}
  \lean{PhyXMiniProblems.ProblemPhyXMini0435.recordedAnswerChoice}
  \uses{def:physics:phyx-mini-0435:phyxminiproblems-problemphyxmini0435-answerchoice}
  Answer label recorded by the supplied dataset
\end{definition}
\begin{proof}
  This is the declared model definition of recorded Answer Choice.
\end{proof}
\begin{definition}[Compatible With Displayed Heat]
  \label{def:physics:phyx-mini-0435:phyxminiproblems-problemphyxmini0435-compatiblewithdisplayedheat}
  \lean{PhyXMiniProblems.ProblemPhyXMini0435.CompatibleWithDisplayedHeat}
  \uses{def:physics:phyx-mini-0435:phyxminiproblems-problemphyxmini0435-answerchoice, def:physics:phyx-mini-0435:phyxminiproblems-problemphyxmini0435-displayedhotreservoirheatinjoules}
  Compatibility with the displayed precision. The exact standard-atmosphere calibration gives 2229150 J, while the source prints 2228000 J; a 2000 J tolerance records the source's coarse/truncated atmospheric conversion
\end{definition}
\begin{proof}
  This is the declared model definition of Compatible With Displayed Heat.
\end{proof}
\begin{definition}[Is Unique Closest Displayed Heat]
  \label{def:physics:phyx-mini-0435:phyxminiproblems-problemphyxmini0435-isuniqueclosestdisplayedheat}
  \lean{PhyXMiniProblems.ProblemPhyXMini0435.IsUniqueClosestDisplayedHeat}
  \uses{def:physics:phyx-mini-0435:phyxminiproblems-problemphyxmini0435-answerchoice, def:physics:phyx-mini-0435:phyxminiproblems-problemphyxmini0435-displayedhotreservoirheatinjoules}
  A choice is the unique closest displayed value to the derived heat
\end{definition}
\begin{proof}
  This is the declared model definition of Is Unique Closest Displayed Heat.
\end{proof}
\begin{lemma}[sample Gas Constant And Unlabeled Temperatures]
  \label{lem:physics:phyx-mini-0435:phyxminiproblems-problemphyxmini0435-samplegasconstantandunlabeledtemperatures}
  \lean{PhyXMiniProblems.ProblemPhyXMini0435.sampleGasConstantAndUnlabeledTemperatures}
  \uses{def:physics:phyx-mini-0435:phyxminiproblems-problemphyxmini0435-reversedrectangularidealgascycle, def:physics:phyx-mini-0435:phyxminiproblems-problemphyxmini0435-gastemperatureinkelvin, def:physics:phyx-mini-0435:phyxminiproblems-problemphyxmini0435-matchesproblemandprimaryfigure, def:physics:phyx-mini-0435:phyxminiproblems-problemphyxmini0435-hasphysicalcycleparameters, def:physics:phyx-mini-0435:phyxminiproblems-problemphyxmini0435-satisfiesidealgaslaws}
  The two labeled isotherms and the ideal-gas law determine the sample's nR readout and the two unlabeled vertex temperatures
\end{lemma}
\begin{proof}
  The conclusion follows from the governing relations and figure readouts named in the statement of sample Gas Constant And Unlabeled Temperatures.
\end{proof}
\begin{lemma}[reversed Cycle Leg Energy Readouts]
  \label{lem:physics:phyx-mini-0435:phyxminiproblems-problemphyxmini0435-reversedcyclelegenergyreadouts}
  \lean{PhyXMiniProblems.ProblemPhyXMini0435.reversedCycleLegEnergyReadouts}
  \uses{def:physics:phyx-mini-0435:phyxminiproblems-problemphyxmini0435-reversedrectangularidealgascycle, def:physics:phyx-mini-0435:phyxminiproblems-problemphyxmini0435-workdonebygasinjoules, def:physics:phyx-mini-0435:phyxminiproblems-problemphyxmini0435-heattransferredintogasinjoules, def:physics:phyx-mini-0435:phyxminiproblems-problemphyxmini0435-matchesproblemandprimaryfigure, def:physics:phyx-mini-0435:phyxminiproblems-problemphyxmini0435-usesidealgascvoverr2model, def:physics:phyx-mini-0435:phyxminiproblems-problemphyxmini0435-hasphysicalcycleparameters, def:physics:phyx-mini-0435:phyxminiproblems-problemphyxmini0435-satisfiesidealgaslaws}
  The governing laws determine each signed work and heat transfer. These values are conclusions and do not occur in any premise structure
\end{lemma}
\begin{proof}
  The conclusion follows from the governing relations and figure readouts named in the statement of reversed Cycle Leg Energy Readouts.
\end{proof}
\begin{lemma}[reversed Cycle Heat And Work Readouts]
  \label{lem:physics:phyx-mini-0435:phyxminiproblems-problemphyxmini0435-reversedcycleheatandworkreadouts}
  \lean{PhyXMiniProblems.ProblemPhyXMini0435.reversedCycleHeatAndWorkReadouts}
  \uses{def:physics:phyx-mini-0435:phyxminiproblems-problemphyxmini0435-reversedrectangularidealgascycle, def:physics:phyx-mini-0435:phyxminiproblems-problemphyxmini0435-matchesproblemandprimaryfigure, def:physics:phyx-mini-0435:phyxminiproblems-problemphyxmini0435-usesidealgascvoverr2model, def:physics:phyx-mini-0435:phyxminiproblems-problemphyxmini0435-hasphysicalcycleparameters, def:physics:phyx-mini-0435:phyxminiproblems-problemphyxmini0435-satisfiesidealgaslaws, def:physics:phyx-mini-0435:phyxminiproblems-problemphyxmini0435-hotreservoirheatinjoules, def:physics:phyx-mini-0435:phyxminiproblems-problemphyxmini0435-coldreservoirheatinjoules, def:physics:phyx-mini-0435:phyxminiproblems-problemphyxmini0435-cycleworkinputinjoules}
  Exact SI energy balance for the reversed cycle: Q\_H = Q\_C + W\_in
\end{lemma}
\begin{proof}
  The conclusion follows from the governing relations and figure readouts named in the statement of reversed Cycle Heat And Work Readouts.
\end{proof}
% --- Archon declaration topology end ---
% --- Archon physics formalization source end ---
